\par\noindent\textbf{(1)} 自由传播矩阵 \(P(L)=\begin{pmatrix}1&L\\0&1\end{pmatrix}\)，薄透镜 \(F(f)=\begin{pmatrix}1&0\\-1/f&1\end{pmatrix}\)。任意组合通过顺序矩阵乘法得到 \(M=\begin{pmatrix}A&B\\C&D\end{pmatrix}\)，且 \(\det M=AD-BC=1\)，与理想光线光学的辛结构相符。
\par\noindent\textbf{(2)} 往返本征值满足 \(\lambda^2-(A+D)\lambda+1=0\)。有界光线振荡需特征值在单位圆上，因此 \(|(A+D)/2|<1\) 为严格稳定条件；等号特征值重合属于边缘情况，是否长期有界还取决于Jordan结构，不能直接判定稳定。
\par\noindent\textbf{(3)} 两镜半径 \(R_1,R_2\)、间距 \(L\)，定义 \(g_i=1-L/R_i\)，标准两镜往返矩阵有半迹 \((A+D)/2=2g_1g_2-1\)。因此腔稳定条件 \(0<g_1g_2<1\)；平凹腔一镜平面 \(R_1=\infty\) 即 \(g_1=1\)，另一镜 \(R_2=R\) 时要求 \(0<L<R\)。
\par\noindent\textbf{(4)} Gaussian束复参数满足光线矩阵的Möbius变换 \(q'=(Aq+B)/(Cq+D)\)。稳定腔自洽模求固定点 \(Cq^2+(D-A)q-B=0\)，取 \(\operatorname{Im}q>0\) 的物理解；几何稳定边界通常使有限束腰解退化，不能仅凭线性光线稳定保证无限高数值孔径或无衍射损耗。
